Este guia apresentou de forma sistemática os fundamentos conceituais e operacionais do pareamento, posicionando-o dentro do arcabouço geral de inferência causal discutido nos outros manuais da série. Assim como no caso de estratégias experimentais e quase-experimentais, o objetivo central permanece o mesmo: recuperar um contrafactual plausível que permita estimar efeitos causais com rigor e transparência.

O pareamento oferece uma resposta particularmente útil para contextos em que a aleatorização não é possível e em que a construção de um grupo de controle comparável depende necessariamente das características observáveis. Ao aproximar a distribuição dessas covariadas entre tratados e controles, minimizamos o viés de seleção decorrente da alocação não aleatória do tratamento e nos aproximamos, tanto quanto os dados permitem, da lógica experimental discutida nos outros guias.

A validade do estimador repousa sobre hipóteses substantivas e a qualidade da inferência depende não apenas das fórmulas, mas da credibilidade empírica dessas premissas. O pareamento exige atenção ao desenho analítico, especialmente na escolha das métricas de distância, número de vizinhos, calipers e verificações de balanceamento.

O fortalecimento dessa etapa diagnóstica é, portanto, parte essencial da prática responsável de pareamento: não basta aplicar o algoritmo; é preciso demonstrar que o procedimento de fato melhora a comparabilidade entre grupos, reduzindo vieses e preservando variabilidade suficiente para garantir precisão. Tal como no planejamento de experimentos, as decisões tomadas antes da estimação — seleção de covariadas, especificações do \textit{propensity score}, definição da amostra — são tão importantes quanto o estimador final.

Por fim, vale ressaltar que o pareamento não elimina, por si só, problemas decorrentes de variáveis não observáveis, nem substitui métodos baseados em hipóteses mais fortes, como Diferenças-em-Diferenças (DID), quando estes são adequados. Ele é, no entanto, uma ferramenta poderosa e intuitiva para aprimorar comparações em dados observacionais, servindo tanto como estimador principal quanto como componente de estratégias híbridas, como DID com pareamento, weighting baseado em \textit{propensity score}, ou análises de sensibilidade.

Com o domínio dos conceitos apresentados — causalidade, resultados potenciais, viés de seleção, métricas de distância, \textit{propensity score}, diagnóstico e interpretação — o leitor está preparado para implementar análises de pareamento de forma rigorosa e alinhada às melhores práticas metodológicas internacionais. Assim como reforçado nos outros guias, a combinação entre clareza conceitual, atenção às hipóteses e disciplina analítica permanece o elemento central de qualquer avaliação causal robusta.